%% Custom operators
\DeclareMathOperator{\spanop}{Span} % Define \spanop for span operator
\DeclareMathOperator{\nullop}{null}
\DeclareMathOperator{\colsp}{colsp} % Define \colsp for column space
\DeclareMathOperator{\image}{image} % Define \image for image operator
\DeclareMathOperator{\Arg}{Arg} % Define \Arg for argument operator
\DeclareMathOperator{\gm}{gm} % Define \gm for geometric multiplicity
\DeclareMathOperator{\am}{am} % Define \am for algebraic multiplicity
%% /Custom operators

\definecolor{lightmagenta}{rgb}{0.95,0.7,0.95}
\definecolor{lightblue}{rgb}{0.7,0.85,0.95}
\definecolor{lightgreen}{rgb}{0.7,0.95,0.7}

% For bemærk boxes
\newcommand{\bemærk}[1]{%
    \tcbox[
        colframe=lightmagenta,
        boxrule=0.5pt,
        arc=2pt,
        left=6pt,
        right=6pt,
        top=3pt,
        bottom=3pt,
        before upper={\textbf{Bemærk: }},
        on line
    ]
    {#1}%
}

\newcommand{\bemærkBig}[1]{%
    \begin{tcolorbox}[
        colframe=lightmagenta,
        boxrule=0.5pt,
        arc=2pt,
        left=6pt,
        right=6pt,
        top=6pt,
        bottom=6pt,
        before upper={\textbf{Bemærk: }},
        before skip=10pt,    % Space before the box
        after skip=10pt,     % Space after the box
    ]
    #1
    \end{tcolorbox}%
}
% /For bemærk boxes

% For Kilde boxes
\newcommand{\kilde}[1]{%
    \hfill
    \tcbox[
        colframe=lightblue,
        boxrule=0.5pt,
        arc=2pt,
        left=6pt,
        right=6pt,
        top=3pt,
        bottom=3pt,
        on line
    ]
    {\kildeformat#1 \kildeend}%
}

\def\kildeformat#1 #2\kildeend{\textbf{#1} #2}

\newcommand{\kildetag}[1]{%
    \tag*{%
        \tcbox[
            colframe=lightblue,
            boxrule=0.5pt,
            arc=2pt,
            left=2pt,
            right=2pt,
            top=0pt,
            bottom=0pt,
            on line
        ]
        {\formatkildetext{#1}}%
    }%
}

\newcommand{\formatkildetext}[1]{%
    \formatkildetextaux#1 \relax\relax%
}
\def\formatkildetextaux#1 #2\relax{%
    \textbf{#1}%
    \ifx\relax#2\relax\else\ #2\fi%
}
% /For Kilde boxes

% For hidden sections
\newcommand{\hiddensection}[1]{%
    \refstepcounter{section}%
    \sectionmark{#1}%
    \addcontentsline{toc}{section}{\texorpdfstring{\protect\sectionSize #1}{#1}}%
}

\newcommand{\hiddensubsection}[1]{%
    \refstepcounter{subsection}%
    \sectionmark{#1}%
    \addcontentsline{toc}{subsection}{\texorpdfstring{\protect\subsecSize #1}{#1}}%
}
% /For hidden sections

% Change real and imaginary part symbols to 2 letters
\renewcommand{\Re}{\mathfrak{Re}}
\renewcommand{\Im}{\mathfrak{Im}}
% /Change real and imaginary part symbols to 2 letters

% /For Eksempel boxes
\newcommand{\eksempel}[1]{%
    \begin{tcolorbox}[
        colframe=lightgreen,
        colback=white,
        boxrule=1 pt,
        arc=2pt,
        left=6pt,
        right=6pt,
        top=6pt,
        bottom=6pt,
        before upper={\textbf{Eksempel: }},
        before skip=10pt,
        after skip=10pt,
    ]
    #1
    \end{tcolorbox}%
}
% /For Eksempel boxes

% Change ugly empty set symbol
\renewcommand{\emptyset}{\text{Ø}}
% /Change ugly empty set symbol

\newcommand*\cleartoleftpage{%
    \clearpage
    \ifodd\value{page}\hbox{}\newpage\fi
}